\section{Introduction}
\label{sec:introduction}

The thermal behaviour of an exothermic continuous stirred tank reactor (CSTR) is set by the competition between heat generation, which grows exponentially with temperature through the Arrhenius law, and heat removal, which grows at most linearly. Since the seminal work of Aris and Amundson \cite{aris1958analysis} and Uppal, Ray and Poore \cite{uppal1974dynamic} this competition is known to produce multiple steady states, open-loop unstable middle branches, limit cycles and parametric sensitivity. The textbook treatment of reactor control \cite{seborg2010process,bequette2003process,luyben2007chemical,ogata2010modern} adds two practical facts: the coolant valve has finite stroke and finite speed \cite{hippe2006windup,tarbouriech2009antiwindup,bohn1995analysis,peng1996nonlinear}, and the heat-transfer coefficient of the jacket decays during operation because of fouling \cite{bott1995fouling,epstein1983thinking,macchietto2011fouling,coletti2010impact}. In industrial practice the reactor temperature is regulated by a cascade, in which a master loop on the reactor temperature sets the jacket temperature and a slave loop moves the coolant valve \cite{shinskey1996process,seborg2010process,astrom2006pid,visioli2006practical}.

These three ingredients---open-loop instability, actuator limits and loss of cooling capacity---are usually discussed together under the label ``thermal runaway''. They are, however, different mechanisms with different remedies. An open-loop unstable operating point can be stabilized by feedback if the controller is fast enough and the actuator has authority; a saturated actuator cannot be helped by any control law; and a degraded heat exchanger changes the plant seen by a fixed-gain controller long before it removes the steady state. A controllability study must therefore say which of these limits is active, and where.

This paper makes that separation on a well-documented three-state benchmark: a first-order irreversible exothermic reaction in a CSTR with a dynamic cooling jacket, the same model used in our companion studies of robust PID tuning \cite{andrade2026robust} and nonlinear predictive control \cite{andrade2026when}. Three questions are addressed:
\begin{enumerate}
    \item Does the benchmark actually have an open-loop unstable steady state, and where? We find that with the nominal feed temperature it does not; steady-state multiplicity appears only when the feed is cooled. All results on the unstable saddle are therefore obtained at a cooled feed, $T_f=\TfSad$~K, and stated as such.
    \item How should a single-loop and a cascade controller be tuned on the saddle, and why does the cascade work? We combine the simple internal model control (SIMC) rules \cite{skogestad2003simple,grimholt2018optimal} with an exact analysis of the closed-inner-loop dynamics, and show that the open-loop gain ratio commonly used for the master design can have the wrong sign at an unstable point.
    \item What happens when the heat-transfer coefficient decays continuously? We introduce a fouling law driven by the reaction rate and compare the time at which the valve saturates with a closed-form quasi-steady prediction. At the stable nominal point the prediction holds; at the saddle the fixed-gain loops are lost much earlier, for a reason that has nothing to do with the cooling capacity.
\end{enumerate}

The contributions are the following. (i) A map of the steady-state branches of the benchmark versus coolant flow and feed temperature, with a closed-form expression for the critical heat-transfer capacity below which a temperature setpoint cannot be held with the valve fully open. (ii) A derivation showing that the condition $\tau_c<\tau_u-\theta$ for an unstable first-order-plus-delay model---a form of the requirement that the closed loop be faster than the unstable pole \cite{stein2003respect,skogestad2005multivariable}---is exact for proportional control of this reactor but not sufficient once the SIMC integral time is added. (iii) The exact static gain of the master loop from the concentration--temperature sub-block of the Jacobian, and its comparison with the open-loop ratio approximation. (iv) A reaction-driven fouling model, a time-to-saturation formula, and the observation that fouling destabilizes fixed-gain loops on the saddle because it shrinks the actuator gain of the coolant valve. All numbers in the paper are generated by one script and are traceable to CSV files.

Section~\ref{sec:model} presents the model, the multiplicity analysis and the thermodynamic limit. Section~\ref{sec:design} derives the controllers. Section~\ref{sec:fouling} introduces the fouling model. Section~\ref{sec:results} reports the simulations and Section~\ref{sec:conclusion} concludes.
